\lecture{7}{Fri 16 Sep 2022 10:31}{}

\begin{theorem}(2.6.2)
	Suppose $N \triangleleft G$. Then there exists a surjective homomorphism $\psi: G \to  G / N$, such that $\operatorname{ker} \psi = N$.
\end{theorem}
\begin{proof}
	We define the \textit{canonical} mapping:
	\begin{align*}
		\psi: G &\longrightarrow G / N \\
		a &\longmapsto \psi(a) = N a
	.\end{align*}
\	Properties to check:
	\begin{itemize}
		\item (Homomorphism) For all $a,b \in  G$, \[
			\psi(ab) = Nab = NaNb = \psi(a)\psi(b)
		.\] 
		\item (surjective) Given any $Na \in G / N$, we have $\psi(a) = Na$.
		\item (kernel) $\operatorname{ker} \psi = \{a \in G: \psi(a) = Ne = N\} $. But $\psi(a) = Na = N$ if and only of $a \in N$. Thus $\operatorname{ker} \psi = N$. 
	\end{itemize}
\end{proof}

\begin{definition}
	If $G$ is a (finite) group and $H \le G$, the number of right cosets of $H$ in $G$ is called the \textit{index} of $H$ in $G$, written $i_G(H)$.
\end{definition}
\begin{notation}
	Also written as $[G:H]$
\end{notation}

\begin{theorem}(2.6.3)
	When $G$ is a finite group and $N \triangleleft G$, then one has \[
		\left| G / N \right| = |G| / |N|
	.\] 
\end{theorem}
\begin{proof}
	Use the proof of Lagrange's Theorem: $|G / N|$ is equal to the right cosets of $N$ in $G$, which is $|G| / |N|$. 
\end{proof}

\begin{observation}
	A group may have no non-trivial normal subgroups. Example: $(\mathbb{Z}_p, +)$ for $p$ prime. 

\end{observation}
\begin{definition}[Simple Groups]
	A group with no non-trivial normal subgroups is called \textit{simple}.
\end{definition}
\begin{remark}
	The classification of simple groups.
\end{remark}

\begin{theorem}[Cauchy's Theorem for Abelian Groups]
	(2.6.4)
	Suppose $G$ is an abelian finite group of order $|G|$, and $p \mid |G|$ with $p$ prime. Then $G$ has an element of order $p$.
\end{theorem}
\begin{proof}
	Proceed by induction on $n = |G|$. The conclusion if trivial for $n = 1$. We suppose $n > 1$ and that the conclusion holds for \textbf{all} abelian groups $G$ with order $1\le |G|<n$. We now consider an abelian group $G$ of order $n$.

	Possibly, we have that $G$ has no non-trivial subgroups. In this case, $G = \left\langle a \right\rangle $ for some $a \in G\setminus \{e\} $, and $|G|$ is prime. Then the order of $a$ is equal to $|G|$, and $G$ has an element of prime order.

	If $G$ has a non-trivial subgroup $N$. Since $G$ is abelian, one has $N \triangleleft G$. Then \[
		|G / N| \cdot |N| = |G| \equiv 0 \pmod p
	.\] 
	So $p \mid |N|$ or $p \mid |G / N|$. In the first case we are done by using induction hypothesis. 

	At this point we can assume that $N \triangleleft  G$, $1 < |N| < |G|$, and $p \nmid |N|$. But $|G / N| |N| = |G|$ and $p \mid  |G|$, we have $p \mid |G / N|$. By the induction hypothesis, since $1 < |G / N|< |G|$, the group  $G / N$ contains an element of order $p$, say $Na$.
	Suppose that $a$ has order $m$ in $G$. Then \[
		(Na)^m = N a^m = Ne = N
	,\] 
	so since $Na$ has order $p$, we must have $p | m$. We write $m  = k p$ with $ k \in \mathbb{N}$, and observe that $a^k$ has order $p$ in $G$.

	So $G$ does indeed have an element of order $p$, completing the induction.
\end{proof}

